\section{Ethics, safety, and reproducibility}
\label{sec:professional_issues}

This work evaluates systems behaviour on public model weights; it does not collect or analyse user conversations. The following points matter when deploying the open-source \texttt{tensor\_store} crate in production.

\textbf{Provenance and safety.} Checkpoints from public hubs should be loaded with the same care as any third-party binary: verify distributor checksums when available, pin revisions, and prefer safe formats (e.g.\ SafeTensors) over unsafe deserialisation. Kernel-adjacent code (\texttt{io\_uring}, unsafe Rust) should use pinned dependency versions and regression tests on upgrades.

\textbf{Privacy and regulation.} Benchmarks do not process tenant data; production deployments must still meet access control, encryption, and data-protection requirements where personal data or proprietary weights are involved.

\textbf{Environmental proportionality.} Faster cold loads can reduce wasted idle time but do not address the dominant energy cost of large-model serving (inference and memory traffic). Claims should separate ``faster loading'' from ``lower total footprint'' unless measured end-to-end.

\textbf{Transparency.} Performance claims should state cache state, format, and layer (Rust, Python, or vLLM). Negative rows (e.g.\ where \texttt{default} trails explicit sync) are reported in the main matrices and limitations. The repository root \texttt{README.md} suggests commands to record kernel, toolchain, and optional GPU and vLLM lines beside result tables. The name \emph{Tensora} refers to this work's system (\texttt{tensor\_store}); it is not affiliated with Google's unrelated TensorStore project for N-dimensional arrays.

\textbf{Integrity.} Upstream engines (vLLM), formats, and community tooling should be cited when results depend on them. Use of generative tools for prose or code should follow institutional policies.

These points align with common professional codes for computing practitioners~\cite{bcs_code}.
